\section{The life, the witnesses, and the question}\label{sec:aug-method}

Augustine of Hippo (354--430) is unusually knowable for
an ancient person. He wrote about childhood, ambition, friendship, desire,
intellectual error, conversion, pastoral burdens, controversy, aging, and the
revision of his own books. A friend and fellow bishop, Possidius of Calama,
survived him and wrote a life that preserves the shape of Augustine's public
ministry and final illness. Augustine's sermons, letters, and treatises place
him in the institutions and conflicts of Roman Africa; conciliar records and
opponents' fragments would add controls but were not comprehensively collated
for this revision. Even the checked abundance creates its own danger. The
\emph{Confessions} is a prayer composed after Augustine became bishop and well
after the Milanese conversion it narrates; the conventional c.~397--401 range
lacks a checked modern chronological locus in this revision. It is neither a diary nor a court
transcript. Possidius writes an exemplary episcopal life, emphasizing Catholic
unity, virtue, and a holy death. Augustine's polemics state real convictions
but present opponents through controversy.

This biography therefore asks not merely, “What happened?” but, “What kind
of witness says that it happened, when, and for what purpose?” It treats
Augustine as a Roman African Christian, rhetor and philosopher, monk, priest,
bishop, Father, and Doctor of the Church. Its geographical center is late Roman
North Africa, with necessary journeys to Italy; its linguistic center is
Latin, though Greek, Punic, and the multilingual world of the Mediterranean
matter. It follows the life chronologically, embeds the writings in the
pressures that produced them, and gives sustained attention to the coercion of
Donatists, anti-Pelagian controversy, gender and patronage, and the difference
between Augustine himself and later “Augustinian” institutions.

\subsection{Evidence key}

\begin{evidencelegend}
\evidencelegendrow{A}{Subject-authored evidence: Augustine's books, letters, sermons, and retrospective self-account, interpreted according to genre and date.}
\evidencelegendrow{N}{Near-contemporary external witness: especially Possidius and, when separately collated, proceedings, laws, and correspondence not controlled by Augustine alone.}
\evidencelegendrow{E}{Early received tradition: late-antique and early-medieval memory, including received death-date, cult, and relic traditions whose individual witnesses still require their own verification.}
\evidencelegendrow{L}{Later hagiographic or devotional tradition: legends, relic narratives, iconographic conventions, and institutional memory.}
\evidencelegendrow{M}{Official ecclesial reception: liturgical, papal, conciliar, or institutional acts, limited to what the act establishes.}
\evidencelegendrow{K}{Modern critical reconstruction: a named scholarly conclusion, including alternatives and degree of confidence.}
\evidencelegendrow{S}{Source-grounded project synthesis: an editorial connection supported by named evidence rather than attributed to one witness.}
\end{evidencelegend}

These codes describe kinds of evidence, not a ladder from unbelief to faith.
Augustine's report of grace in his conversion is the testimony of a Christian
subject to God's action. Historical criticism can locate the report, its
literary design, and its relationship to earlier evidence; it cannot repeat
the experience as an experiment or disprove it by relabeling it psychology.
Likewise, Possidius's accounts of healing are reported as Possidius reports
them (\emph{Life} 29), without this study independently certifying a miracle or
inventing a natural explanation.

\subsection{Chronological orientation}

\begin{biotimeline}
\biotimelinerow{354}{Thagaste, Numidia}{Born to Patricius and Monnica on 13 November; Augustine gives the birthday in \emph{De beata vita} 1.6, while the year is secured by the converging ancient chronology.}{Well attested}
\biotimelinerow{c.~365--370}{Madaura / Thagaste}{Grammar and literary education; a financially constrained interval at home precedes advanced study. \emph{Confessions} 1--2 is the retrospective source.}{Approximate}
\biotimelinerow{371--383}{Carthage and Africa}{Rhetorical studies, a long partnership, birth of Adeodatus, teaching, adherence to Manichaeism, and growing dissatisfaction; Cicero's lost \emph{Hortensius} redirects his desire for wisdom (\emph{Confessions} 3--5).}{Sequence secure; dates partly approximate}
\biotimelinerow{383--386}{Rome / Milan}{Moves to Rome, then receives the imperial rhetoric chair at Milan; hears Ambrose, encounters Platonist books, and abandons a court marriage project after his partner's dismissal.}{Well attested in A}
\biotimelinerow{386--387}{Milan / Cassiciacum}{Conversion crisis in late summer 386; withdrawal to Cassiciacum; baptism with Adeodatus and Alypius at the Easter Vigil of 387.}{Event secure; exact conversion day reconstructed}
\biotimelinerow{After 387}{Ostia / Rome / Africa}{Monnica dies at Ostia; Augustine eventually returns to Africa and establishes a lay ascetic community at Thagaste.}{Relative sequence secure; return date unresolved}
\biotimelinerow{391}{Hippo Regius}{Ordained presbyter by Bishop Valerius after visiting Hippo; establishes a monastery and begins sustained preaching and controversy.}{Well attested}
\biotimelinerow{395/396}{Hippo Regius}{Consecrated coadjutor and succeeds Valerius as bishop.}{Year boundary disputed}
\biotimelinerow{After episcopal succession--411}{Roman Africa}{Composes the \emph{Confessions}; expands pastoral and anti-Donatist work; participates in the 411 Conference of Carthage.}{Relative sequence secure; work dates partly unresolved}
\biotimelinerow{Later episcopate, after 410}{Hippo and wider networks}{Anti-Pelagian writing, long works on the Trinity and the two cities, extensive correspondence, and preaching.}{Relative sequence; detailed work chronology uncollated}
\biotimelinerow{426}{Hippo}{Designates the presbyter Heraclius to assume pastoral burdens; Letter 213 records the event. Late in his episcopate he begins reviewing his books in the \emph{Retractationes}.}{Heraclius dated; self-review start year unresolved}
\biotimelinerow{430}{Hippo Regius}{Dies during the Vandal siege, aged seventy-five by modern reckoning, and is buried locally. Possidius gives the siege context, age, death, and burial.}{Death secure; the received 28 August date is not independently collated here}
\end{biotimeline}

\begin{studyframe}
\textbf{Governing thesis.} Augustine's sanctity cannot responsibly be reduced
either to a private conversion in a garden or to an abstract system called
“Augustinianism.” It took historical form in a life of prayer, intellectual
labor, common life, preaching, sacramental ministry, institutional power,
friendship, correction, and controversy. The same record that reveals unusual
self-criticism also requires judgment about his defenses of religious
coercion and the harsh social afterlife of some formulations.
\end{studyframe}
